\documentclass[12pt]{article}
\usepackage{amsmath, amssymb, amsthm}
\usepackage{parskip}
\usepackage{geometry}
\usepackage{enumitem}
\usepackage{fancyhdr}

\geometry{letterpaper, margin=1in}
\setlength{\headheight}{15pt}

% Header setup
\pagestyle{fancy}
\fancyhf{}
\lhead{Math 311 Fall 2026}
\rhead{Jeffrey Lu}
\cfoot{\thepage}

\title{Homework Assignment 6 \\ \large Math 311, Fall 2026 \\ Rutgers University}
\author{}
\date{}

\begin{document}

\maketitle

\textbf{Exercise 3.2.2}
Let
\[
    A = \{ (-1)^n + \frac{2}{n} : n = 1, 2, 3, \dots \} \quad \text{ and } \quad  B = \{x \in \mathbb{Q} : 0 < x < 1 \}.
\]
\begin{enumerate}[label=(\alph*)]
    \item For set \( A \), we can seperate into even and odd indices.
          If \( n \) is even
          \[
              (-1)^{2n} + \frac{2}{2n} = 1 + \frac{1}{n} \to 1.
          \]
          If \( n \) is odd
          \[
              (-1)^{2n-1} + \frac{2}{2n-1} = -1 + \frac{2}{2n-1} \to -1.
          \]
          Thus, the limit points of \( A \) are \( -1, 1 \).

          For  \( B \), since \( \mathbb{Q} \) is dense in \( \mathbb{R} \)
          all elements in \( [0,1] \) are limit points of \( B \).

    \item The set \( A \) is not closed since \( -1 \) is a limit point of \( A \) but \( -1 \not \in A \).
          The set \( A \) is also not open since \( 2 \in A \), but every \( \epsilon \)-neighborhood of \( 2 \) contains elements not in \( A \).
          Therefore, \( A \) is neither open or closed.

          The set \( B \) is not closed, since \( 0,1 \not \in B \) but both are limit points.
          The set \( B \) is not open, since for each \( x \in B \), its neighborhood contains irrationals.
          Therefore, \( B \) is neither open or closed.

    \item Every element of \( A \setminus \{1\} \) is an isolated point.

          The set \( B \) has no isolated points due to the density of \( \mathbb{Q} \).

    \item The closure of \( A \) is
          \[
              \overline{A} = A \cup \{-1\}.
          \]
          The closure of \( B \) is
          \[
              \overline{B} = [0,1].
          \]
\end{enumerate}

\textbf{Exercise 3.2.4}
Let \( A \) be nonempty and bounded above so that \( s = \sup A \) exists.
\begin{enumerate}[label=(\alph*)]
    \item We proceed by splitting into two cases.
          If \( s \in A \), then \( s \in \overline{A}  \) by definition of closure.

          If \( s \not \in A \), let \( \epsilon > 0 \).
          Since \( s = \sup A \), we have \( a \in A \) such that
          \[
              s - \epsilon < a < s.
          \]
          So
          \[
              0 < \vert a - s \vert < \epsilon.
          \]
          Hence
          \[
              a \in V_{\epsilon}(s) \cap (A \setminus \{s\})
          \]
          and \( s  \) is a limit point of \( A \).
          Therefore, in both cases \( s \in \overline{A} \).

    \item No, suppose, for contradiction, that \( A \) is an open set and \( s = \sup A \in A \).
          Since \( A  \) is open, there exists \( \epsilon > 0 \) such that
          \[
              V_{\epsilon}(s) = ( s - \epsilon, s + \epsilon) \subseteq A.
          \]
          So \( s + \frac{\epsilon}{2} \in A \). \
          But \( s + \frac{\epsilon}{2} > s \), contradicting the fact that \( s \) is an upper bound.
          Therefore, \( s \not \in A \).


\end{enumerate}

\textbf{Exercise 3.2.6}
\begin{enumerate}[label=(\alph*)]
    \item False, consider \( A = \mathbb{R} \setminus \{\sqrt{2}\} \).
          Since \( \{\sqrt{2}\} \) is a closed set, its complement \( A \) must be open.
          But \( A \neq \mathbb{R} \) even though it is an open set containing every rational.

    \item False, consider \( F_n = [n, \infty ) \).
          Then for each \( n \in \mathbb{N} \), \( F_n \) is closed and nonempty, and
          \[
              F_1 \supseteq  F_2 \supseteq  F_3 \supseteq  \dots
          \]
          But
          \[
              \bigcap_{n=1}^{\infty} F_n = \varnothing.
          \]

    \item True, let \( A \subseteq \mathbb{R} \) be a nonempty open set.
          Choose some \( x \in A \).
          Since \(A \) is open, there exists \( \epsilon > 0 \) such that
          \[
              V_{\epsilon}(x) = (x-\epsilon, x+\epsilon) \subseteq A.
          \]
          By density of \( \mathbb{Q} \), there exists \( q \in \mathbb{Q} \) such that
          \[
              x - \epsilon < q < x + \epsilon.
          \]
          Hence
          \[
              q \in V_{\epsilon}(x) \subseteq A.
          \]
          Therefore, \( A \) must contain a rational number.


    \item False, consider
          \[
              A = \{\sqrt{2} + \frac{1}{n} : n \in \mathbb{N} \} \cup  \{\sqrt{2}\}.
          \]
          Then \( A \) is a infinite, bounded between \( [\sqrt{2}, \sqrt{2} + 1] \),
          closed since it contains its limit point \( \sqrt{2} \), and contains no rationals
          since the sum of a irrational and a rational is irrational.


    \item True, recall that the cantor set starts with \( [0,1] \) and we repeatedly
          remove the open middle.
          This leaves us with a finite union of closed intervals, which we know is closed.
          Then we take the intersection of these closed intervals, which is also closed by theorem 3.2.14

\end{enumerate}

\textbf{Exercise 3.2.9}
\begin{enumerate}[label=(\alph*)]
    \item For
          \[
              \left( \bigcup_{\lambda \in \Lambda } E_{\lambda } \right)^c = \bigcap_{\lambda \in \Lambda} E_{\lambda }^c
          \]
          Let \( x \in \left( \bigcup_{\lambda \in \Lambda } E_{\lambda } \right)^c \). Then
          \[
              \begin{split}
                  x \in \left( \bigcup_{\lambda \in \Lambda } E_{\lambda } \right)^c
                   & \iff x \not \in \bigcup_{\lambda \in \Lambda } E_{\lambda}              \\
                   & \iff x \not \in E_{\lambda} \quad \text{for every } \lambda \in \Lambda \\
                   & \iff x \in E_{\lambda}^c \quad \text{for every } \lambda \in \Lambda    \\
                   & \iff x \in \bigcap_{\lambda \in \Lambda } E_{\lambda}^c
              \end{split}
          \]
          Therefore,
          \[
              \left( \bigcup_{\lambda \in \Lambda } E_{\lambda } \right)^c = \bigcap_{\lambda \in \Lambda} E_{\lambda }^c
          \]

          For
          \[
              \left(\bigcap_{\lambda\in\Lambda}E_\lambda\right)^c
              = \bigcup_{\lambda\in\Lambda}E_\lambda^c.
          \]
          Let \( x \in \left(\bigcap_{\lambda\in\Lambda}E_\lambda\right)^c \).
          Then
          \[
              \begin{split}
                  x \in \left(\bigcap_{\lambda\in\Lambda}E_\lambda\right)^c
                   & \iff x \not \in \bigcap_{\lambda\in\Lambda}E_\lambda                   \\
                   & \iff x \not \in E_{\lambda} \quad \text{for some } \lambda \in \Lambda \\
                   & \iff x \in E_{\lambda}^c \quad \text{for some } \lambda \in \Lambda    \\
                   & \iff x \in \bigcup_{\lambda\in\Lambda}E_\lambda^c.
              \end{split}
          \]
          Therefore,
          \[
              \left(\bigcap_{\lambda\in\Lambda}E_\lambda\right)^c
              = \bigcup_{\lambda\in\Lambda}E_\lambda^c.
          \]


    \item Let \( \{E_{\lambda} : \lambda \in \Lambda\} \) be a finite collection of closed sets.
          By theorem 3.2.13, for each \( \lambda \in \Lambda \), \( E_\lambda^c \) is open.
          Then by theorem 3.2.3, the finite intersection of these open sets is also open.
          Since the complement of this intersection is closed, it follows that
          \[
              \left(\bigcap_{\lambda\in\Lambda}E_\lambda^c\right)^c = \bigcup_{\lambda\in\Lambda}E_\lambda.
          \]
          Hence the union of a finite collection of closed set is closed.


          Let \( \{F_{\lambda} : \lambda \in \Lambda\} \) be a arbitrary collection of closed sets.
          By theorem 3.2.13, for each \( \lambda \in \Lambda \), \( F_\lambda^c\) is open.
          Then by theorem 3.2.3, the union of these open sets is also open.
          Since the complement of this union must be closed, it follows that
          \[
              \left( \bigcup_{\lambda \in \Lambda } F_{\lambda}^c \right)^c = \bigcap_{\lambda \in \Lambda} F_{\lambda }
          \]
          Hence the intersection of an arbitrary collection of closed sets is closed.
\end{enumerate}

\textbf{Exercise 3.2.15}
\begin{enumerate}[label=(\alph*)]
    \item For each \( n \in \mathbb{N} \), define
          \[
              G_n = (a - \frac{1}{n}, b + \frac{1}{n}).
          \]
          Then each \( G_n \) is an open set.

          First, let \( x \in [a,b] \). Then for every \( n \in \mathbb{N} \)
          \[
              a - \frac{1}{n} < x < b + \frac{1}{n}.
          \]
          Thus \( x \in G_n \) for every \( n \in \mathbb{N} \), so
          \[
              [a,b] \subseteq \bigcap_{n=1}^{\infty} G_n.
          \]

          Let
          \[
              x \in \bigcap_{n=1}^{\infty} G_n.
          \]
          Then \( x \in G_n \) for every \( n \in \mathbb{N} \).
          So
          \[
              a - \frac{1}{n} < x < b + \frac{1}{n}.
          \]
          Since \( \frac{1}{n} \to 0 \), by the Order Limit Theorem
          \[
              a \leq x \leq b.
          \]
          Hence \( x \in [a,b] \), and so
          \[
              \bigcap_{n=1}^{\infty} G_n \subseteq [a,b].
          \]
          Therefore we have
          \[
              \bigcap_{n=1}^{\infty} G_n = [a,b].
          \]


    \item For each \( n \in \mathbb{N} \) we define
          \[
              G_n = (a, b + \frac{1}{n}).
          \]
          Then each \( G_n \) is an open set.
          Let \( x \in (a, b] \). Then
          \[
              a < x \leq b < b + \frac{1}{n}.
          \]
          for every \( n \in \mathbb{N} \).
          So \( x \in G_n \) for every \( n \in \mathbb{N} \), and
          \[
              x \in \bigcap_{n=1}^{\infty} G_n.
          \]
          Hence
          \[
              (a,b] \subseteq \bigcap_{n=1}^{\infty} G_n.
          \]

          Let
          \[
              x \in \bigcap_{n=1}^{\infty} G_n.
          \]
          Then \( x \in G_n \) for every \( n \in \mathbb{N} \).
          So
          \[
              a < x < b + \frac{1}{n}.
          \]
          Since \( \frac{1}{n} \to 0 \), the Order Limit theorem says
          \[
              a < x \leq b.
          \]
          So \( x \in (a,b] \), and
          \[
              \bigcap_{n=1}^{\infty} G_n \subseteq (a,b].
          \]
          Therefore
          \[
              (a,b] = \bigcap_{n=1}^{\infty} G_n.
          \]

          Next, we show \( (a,b] \) is an \( F_{\sigma} \) set.
          For each \( n \in \mathbb{N} \) we define
          \[
              F_n = [a + \frac{1}{n}, b]
          \]
          where \( F_n = \varnothing  \) if \( a +\frac{1}{n} > b \).
          Then each \( F_n \) is a closed set.
          Let \( x \in (a,b] \).
          Then
          \[
              a \leq x \leq b.
          \]
          Since \( x - a > 0 \), by the Archimedean Property, there exists \( n \in \mathbb{N} \)
          such that
          \[
              \frac{1}{n} < x - a.
          \]
          Hence
          \[
              a + \frac{1}{n} < x \leq b.
          \]
          That is, \( x \in F_n \) for some \( n \in \mathbb{N} \). So
          \[
              x \in \bigcup_{n=1}^{\infty} F_n.
          \]
          Let
          \[
              x \in \bigcup_{n=1}^{\infty} F_n.
          \]
          Then there exists some \( n \in \mathbb{N} \) such that \( x \in F_n \).
          For this \( n \), we have
          \[
              a + \frac{1}{n} \leq x \leq b.
          \]
          Sicne \( \frac{1}{n} > 0 \), we have
          \[
              a < x \leq b.
          \]
          Thus, \( x \in (a,b] \), so
          \[
              \bigcup_{n=1}^{\infty} F_n \subseteq (a,b].
          \]

          Therefore,
          \[
              \bigcup_{n=1}^{\infty} F_n = (a,b].
          \]


    \item To show that \( \mathbb{Q} = F_{\sigma } \), recall that \( \mathbb{Q} \) is countable,
          so we can list its elements as
          \[
              \mathbb{Q} = \{q_1, q_2, q_3, \dots \}.
          \]
          Then every element is a closed set, and so
          \[
              \mathbb{Q} = \bigcup_{n=1}^{\infty} \{q_n\}.
          \]

          Next, we show that the irrationals form a \( G_{\delta } \) set.
          Define
          \[
              I = \mathbb{R} \setminus \mathbb{Q}.
          \]
          Then
          \[
              \begin{split}
                  I & = \mathbb{Q}^c                                   \\
                    & = \left(\bigcup_{n=1}^{\infty} \{q_n\} \right)^c \\
                    & = \bigcap_{n=1}^{\infty} \{q_n\}^c.
              \end{split}
          \]
          Since \( \{q_n\} \) is closed for each \( n \in \mathbb{N} \),
          each complement must then be open.
          Therefore, the irrationals are the countable intersection of open sets.
\end{enumerate}

\textbf{Exercise 3.3.1}
Suppose that \( K \) is compact and nonempty.
By Heine-Borel Theorem, \( K \) is also closed and bounded.
Then by the Axiom of Completeness, both \( \sup K \) and \( \inf K \) exists.
Exercise 3.2.4 says that
\( \sup K, \inf K \in \overline{K} \).
Since \( K \) is closed, \( K = \overline{K} \), and so
\[
    \sup K, \inf K \in K.
\]

\textbf{Exercise 3.3.2}
\begin{enumerate}[label=(\alph*)]
    \item Not compact, consider \( x_n = n \).
          For every upper bound \( M \), by Archimedean Property we can always produce \( n > M \).
          Thus \( x_n \) is unbounded and so every subsequence does not converge.

    \item Consider the sequence of the decimal approximation of \( \sqrt{2}/2 \)
          \[
              x_n=0.7,\ 0.70,\ 0.707,\ 0.7071,\ 0.70710,\ldots
          \]
          Then every term is rational belongs \( [0,1] \), but
          \( x_n \to  \sqrt{2}/2 \). Since every subsequence also converges to same point,
          no subsequences converge to a point in \( \mathbb{Q} \cap [0,1] \). Hence it is not compact.


    \item Recall form exercise 3.2.6, we showed that the Cantor Set is closed. Furthermore, since
          \[
              C \subseteq [0,1]
          \]
          , \( C \) is also bounded.
          Thus by Theorem 3.3.4, \( C \) is compact.

    \item Not compact, consider the sequence
          \[
              x_n = \sum_{k=1}^{n} \frac{1}{k^2}.
          \]
          Each \( x_n \in \{1 + \frac{1}{2^2} + \frac{1}{3^2} + \dots \} \), and
          \[
              x_{n+1} - x_{n} = \frac{1}{(n+1)^2} > 0
          \]
          so the sequence is increasing.
          Since for \( k \geq 2 \)
          \[
              \frac{1}{k^2} < \frac{1}{k(k-1)},
          \]
          it follows that \( x_n \) is bounded above.
          Thus,
          \[
              \begin{split}
                  x_n & = 1 + \sum_{k=2}^{n} \frac{1}{k^2}                \\
                      & < 1 + \sum_{k=2}^{n} \frac{1}{k(k-1)}             \\
                      & = 1 + \sum_{k=2}^{n}(\frac{1}{k-1} - \frac{1}{k}) \\
                      & = 2 - \frac{1}{n}                                 \\
                      & < 2.
              \end{split}
          \]
          for \( n \geq 2 \).
          By the Monotone Convergence Theorem, \( (x_n) \to L \in \mathbb{R} \).
          Since \( x_n \) is increasing, by the Order Limit Theorem
          \[
              L \geq x_{n+1} > x_n.
          \]
          Thus, \( L \not \in  \{1 + \frac{1}{2^2} + \frac{1}{3^2} + \dots \}\).
          Since no subsequence converges to a point inside the set, it is not compact.

    \item Since
          \[
              \frac{n}{n+1} = 1 - \frac{1}{n+1} \to 1.
          \]
          Since \( 1 \in B \), \( B \) is closed.
          Also \( B \subseteq [0,1] \), so \( B \) is bounded.
          Thus \( B \) is compact.
\end{enumerate}

\textbf{Exercise 3.3.5}
\begin{enumerate}[label=(\alph*)]
    \item Let \(\{K_\lambda:\lambda\in\Lambda\}\) be an arbitrary collection of compact sets.
          Since each \(K_\lambda\) is compact, each is closed and bounded by Theorem 3.3.4.
          By Theorem 3.2.14, an arbitrary intersection of closed sets is closed. Thus,
          \[
              K=\bigcap_{\lambda\in\Lambda}K_\lambda
          \]

          is closed.
          Moreover, since
          \[
              K\subseteq K_\lambda
          \]

          for every \(\lambda\in\Lambda\), \(K\) is a subset of a bounded set and is therefore bounded (assuming the collection is nonempty).
          Hence \(K\) is closed and bounded, so it is compact by Theorem 3.3.4.

    \item False. For each \(n\in\mathbb N\), let
          \[
              K_n=[-n,n].
          \]

          Every \(K_n\) is closed and bounded, so each is compact.
          However,
          \[
              \bigcup_{n=1}^{\infty}K_n=\mathbb R,
          \]

          which is unbounded and therefore not compact.


    \item False.Let
          \[
              A=(0,1),\qquad K=[0,1].
          \]

          Since \(K\) is closed and bounded, it is compact.
          However,
          \[
              A\cap K=(0,1),
          \]

          which is not closed and therefore not compact.

    \item False. For each \(n\in\mathbb N\), define
          \[
              F_n=[n,\infty).
          \]

          Every \(F_n\) is nonempty and closed. Furthermore,
          \[
              F_1\supseteq F_2\supseteq F_3\supseteq\cdots.
          \]

          However,
          \[
              \bigcap_{n=1}^{\infty}F_n=\emptyset.
          \]

          To see why, suppose there exists \(x\) in the intersection. Then
          \[
              x\ge n\qquad\text{for every }n\in\mathbb N.
          \]

          But by the Archimedean Property, there exists \(n\in\mathbb N\) such that \(n>x\), a contradiction.
          Therefore, the intersection is empty.
\end{enumerate}

\end{document}